\section{The actual native factor on arbitrary original residue vectors}

\begin{theorem}[The native transition factor computes every original residue vector]
\label{thm:pl-original-native-factor-arbitrary-residue-vector}
\lean{EnvelopeProperty.Berkovich.RelativeSchemeValuedPoint.modelLineNorm_centerFrameGenericIso_unitSection,EnvelopeProperty.NativeTestConfiguration.TestConfiguration.gaussCompatibleTargetFrame_residue_unit_norm,EnvelopeProperty.NativeTestConfiguration.TestConfiguration.gaussCompatibleNormFactor_residue_formula}
\leanok
\uses{thm:proper-model-arbitrary-residue-vector-norm}
For an actual test configuration and its original line comparison and
fibre-one identification, choose any original reference frame at the
prescribed original point. The true model norm of every residue vector
at its actual Gauss orbit point equals the constructed native transition
factor times the original reference-frame coefficient value of the
transported residue vector. The vector need not extend to a global
section on the original total space, and the original group product
need not be proper.
\end{theorem}
\begin{proof}
Choose the actual proper-center target frame used in the definition of
the native factor. Its unit section has model norm one, because the
target frame is obtained from a frame on the actual valuation-ring
center. The arbitrary-residue pullback formula expresses the original
model norm as the coefficient value in the induced model frame times
this unit norm. The whole transition unit between the chosen reference
frame and this induced model frame compares the coefficients of every
residue vector, including vectors without global extensions. Its value
is exactly the defined native factor. Substitute the unit norm one
and the actual transition coefficient identity.
\end{proof}
